%%%% CAPÍTULO 3 - MATERIAIS E MÉTODOS

\chapter{Materiais e Métodos}\label{cap:metodologia}

Este capítulo detalha os procedimentos metodológicos adotados para a realização
do estudo, cobrindo as fases de coleta, tratamento e preparação dos dados, bem
como as decisões de projeto que sustentam a reprodutibilidade da análise. A
estrutura da pesquisa fundamenta-se no modelo CRISP-DM, descrito na
Seção~\ref{sec:crisp-dm}.

\section{Ambiente de Desenvolvimento e Tecnologias}\label{sec:ambiente}

O ecossistema de desenvolvimento foi estruturado na linguagem de programação
Python~3. Foram empregadas as seguintes bibliotecas computacionais:

\begin{itemize}
  \item \textbf{pandas:} manipulação de estruturas de dados tabulares,
    tratamento de valores ausentes e transformações matemáticas;
  \item \textbf{NumPy e SciPy:} operações numéricas, testes de hipótese e
    otimização por máxima verossimilhança~\cite{Virtanen2020};
  \item \textbf{scikit-learn:} implementação dos modelos de aprendizado de
    máquina e das rotinas de avaliação~\cite{Pedregosa2011};
  \item \textbf{Matplotlib:} geração das figuras estatísticas.
\end{itemize}

Diferentemente da etapa anterior da pesquisa, na qual a análise foi conduzida
integralmente em um \textit{notebook}, o código foi reorganizado em módulos
Python versionados e cobertos por testes automatizados. Essa reorganização
atende a dois propósitos. O primeiro é a verificabilidade: cada procedimento de
cálculo possui testes que fixam seu comportamento esperado, incluindo casos
construídos para falhar caso um erro específico seja reintroduzido. O segundo é a
separação entre a lógica da análise e a narrativa do trabalho, de modo que o
\textit{notebook} passe a consumir funções testadas em vez de reimplementá-las.

\section{Coleta e Fonte de Dados}\label{sec:coleta}

O conjunto de dados utilizado foi obtido a partir do repositório público
\textit{Football-Data.co.uk}~\cite{FootballData2026}, que disponibiliza o
histórico detalhado das partidas da \textit{Premier League} em arquivos no
formato CSV, organizados por temporada.

O recorte temporal estabelecido abrange 11 temporadas completas, da temporada
2015/16 à temporada 2025/26, totalizando 4.180 partidas. Cada registro contém
variáveis como a identificação dos times mandante e visitante, os gols marcados
por cada equipe, o total de finalizações, as finalizações no alvo, os escanteios
cobrados, as faltas cometidas e os cartões aplicados.

\subsection{Integridade e Procedência dos Dados}\label{sec:integridade}

A fonte reescreve periodicamente seus arquivos, seja para corrigir placares, seja
para acrescentar colunas. Uma coleta realizada diretamente a cada execução, como
na etapa anterior da pesquisa, implica que duas execuções em datas distintas
possam produzir números diferentes sem qualquer indicação disso.

Para evitar esse comportamento, adotou-se o seguinte procedimento: os arquivos
são baixados uma única vez e armazenados localmente; a procedência de cada
arquivo, a data de acesso e seu resumo criptográfico SHA-256 são registrados em
um manifesto; e toda execução subsequente confere os arquivos contra esse
manifesto, interrompendo o processamento em caso de divergência. Falhas de
coleta, por sua vez, interrompem a execução em vez de serem registradas e
ignoradas, o que impede que uma análise prossiga sobre um conjunto incompleto.

\subsection{Verificação de Comparabilidade Histórica}\label{sec:quebras}

Antes de qualquer comparação ao longo do tempo, verificou-se se as próprias
variáveis são comparáveis entre temporadas. Implementou-se, para isso, uma
rotina que sinaliza variações abruptas nas médias da liga entre temporadas
consecutivas.

Aplicada a um recorte estendido, iniciado em 2009/10, a rotina identificou uma
descontinuidade: a média de finalizações no alvo da liga cai aproximadamente
37\% entre as temporadas 2012/13 e 2013/14, enquanto gols e finalizações totais
permanecem estáveis. A estabilidade das demais variáveis indica que não se trata
de mudança no jogo, mas de alteração no critério de registro adotado pelo
provedor dos dados. Em consequência, comparações dessa variável que atravessem a
fronteira de 2013/14 não são válidas. O recorte principal do trabalho situa-se
integralmente após essa descontinuidade e, portanto, não é afetado; a restrição
aplica-se apenas ao estudo de caso apresentado na
Seção~\ref{sec:resultados-estadios}.

\section{Preparação e Engenharia de Dados}\label{sec:preparacao}

Os dados brutos necessitaram de transformações estruturais para que fosse
possível responder aos objetivos específicos da pesquisa.

\subsection{Criação de Variáveis de Desempenho}\label{sec:variaveis}

Foram computadas as seguintes variáveis derivadas:

\begin{itemize}
  \item \textbf{Pontos conquistados:} aplicação da regra oficial do esporte,
    atribuindo 3 pontos para vitória, 1 ponto em caso de empate e 0 pontos para
    derrota, calculados separadamente para mandante e visitante;
  \item \textbf{Saldo de gols do mandante:} diferença entre os gols marcados
    pelo time da casa e os gols sofridos;
  \item \textbf{Gols totais da partida:} soma dos gols de ambas as equipes,
    utilizada como indicador de dinamicidade dos jogos.
\end{itemize}

\subsection{Formato Longo e Perspectiva da Equipe}\label{sec:formato-longo}

O conjunto de dados original organiza-se com uma linha por partida, o que dificulta
agregações por equipe: as métricas de um mesmo clube encontram-se ora nas colunas
de mandante, ora nas de visitante, conforme o jogo. Essa organização favorece um
erro específico, consistente em agregar colunas de mandante e visitante de uma
mesma partida e interpretar o resultado como desempenho dentro e fora de casa de
um mesmo grupo de equipes, quando se trata, na realidade, do desempenho de
grupos distintos.

Para eliminar essa possibilidade, o conjunto foi convertido para o formato longo,
com duas observações por partida — uma para cada equipe — e as métricas sempre
expressas na perspectiva do clube, acompanhadas de uma coluna explícita de mando.
Nesse formato, agregações por clube e por condição de mando tornam-se diretas.

\subsection{Tabela de Classificação}\label{sec:classificacao}

Para possibilitar a análise do comportamento das equipes segundo sua posição
final, construiu-se a tabela de classificação de cada temporada a partir dos
resultados das partidas, ordenada pelo critério oficial da competição: pontos,
saldo de gols e gols marcados. Critérios de desempate subsequentes, como
confronto direto, não foram implementados, por raramente alterarem a faixa de
posição de uma equipe.

As equipes foram então segmentadas em três faixas, recalculadas a cada
temporada: G6, correspondente às posições 1 a 6; Meio, posições 7 a 14; e Z6,
posições 15 a 20. Essa segmentação é preferível a rótulos fixos de grupo, como o
chamado \textit{Big Six}, pois acompanha a variação de desempenho dos clubes ao
longo do período estudado.

\subsection{Classificação do Regime de Público}\label{sec:regime-publico}

A ausência de público não coincide com temporadas inteiras. Apenas 92 das 380
partidas da temporada 2019/20 ocorreram com portões fechados, correspondentes ao
período posterior à retomada da competição, em 17 de junho de 2020. Além disso, a
temporada 2020/21 registrou janelas de liberação parcial de público.

Por esse motivo, a classificação foi realizada por data da partida e não por
temporada, com três categorias: público total, público limitado e sem público. A
categoria intermediária merece justificativa. As autorizações de dezembro de 2020
e maio de 2021 variavam por clube e por semana, e o público efetivo de cada
partida não está disponível na fonte. Atribuir essas partidas a um dos regimes
extremos exigiria suposição não respaldada pelos dados; mantê-las em categoria
própria, excluída dos contrastes, preserva a validade das comparações.

\section{Procedimentos de Análise Estatística}\label{sec:procedimentos}

Os testes de hipótese foram selecionados de modo que a hipótese nula correspondesse
exatamente à pergunta de pesquisa.

Para avaliar a vantagem do mandante, adotou-se o teste binomial condicionado aos
jogos decididos, cuja hipótese nula estabelece igual probabilidade de vitória para
mandante e visitante. Essa formulação é preferível ao teste qui-quadrado contra a
distribuição uniforme entre os três resultados possíveis, pois esta última embute
a suposição de que um terço das partidas terminaria empatada, proporção que
nenhuma teoria do Fator Casa prevê e que não se verifica historicamente.

Para a comparação entre pontos do mandante e do visitante, empregou-se o teste
$t$ pareado. As duas séries provêm da mesma partida e são determinísticas entre
si, de modo que a hipótese de independência exigida pelo teste para amostras
independentes não se sustenta.

Todos os resultados inferenciais são reportados com estimativa pontual, intervalo
de confiança de 95\% e tamanho de efeito, conforme discutido na
Seção~\ref{sec:efeito}. Intervalos para estatísticas sem forma fechada foram
obtidos por reamostragem \textit{bootstrap}.

\section{Modelagem}\label{sec:modelagem}

\subsection{Modelo de Poisson Bivariado}\label{sec:met-dixon-coles}

Adotou-se o modelo de \citeonline{DixonColes1997}, especificado na
Equação~\ref{eq:dixon-coles}, estimado por máxima verossimilhança. A
identificabilidade dos parâmetros de ataque e defesa, determinados a menos de
uma constante aditiva, foi assegurada pela imposição de média nula aos
parâmetros de ataque.

O erro padrão do parâmetro de mando foi obtido pela curvatura da verossimilhança
perfilada, isto é, reotimizando os demais parâmetros a cada deslocamento de
$\gamma$. Essa escolha decorre de verificação empírica: a curvatura marginal,
que mantém os demais parâmetros fixos, subestimou a incerteza em aproximadamente
25\% em simulação controlada, em razão da correlação entre $\gamma$ e o termo de
escala~$\mu$.

Para estimar o efeito de uma condição binária sobre o mando, o parâmetro foi
reparametrizado como $\gamma + \delta \cdot \mathbb{1}[\text{condição}]$. Essa
formulação permite estimar o efeito dentro de um mesmo recorte temporal, com um
único parâmetro adicional, viabilizando a análise em subconjuntos nos quais o
ajuste de um modelo completo não seria possível.

A implementação foi validada por recuperação de parâmetro: partidas sintéticas
foram geradas a partir do próprio modelo, com $\gamma$ conhecido, verificando-se
se o procedimento de ajuste o recupera e se a cobertura dos intervalos de
confiança corresponde ao nível nominal.

\subsection{Modelos Preditivos}\label{sec:met-preditivo}

A etapa preditiva teve como alvo o resultado da partida, em suas três categorias.
Em conformidade com o exposto na Seção~\ref{sec:vazamento}, nenhuma estatística
produzida durante a partida foi utilizada como variável explicativa. As variáveis
empregadas utilizam exclusivamente informação disponível antes do apito inicial:
\textit{rating} Elo de cada equipe, desempenho nas cinco partidas anteriores em
pontos e gols, desempenho recente específico da condição de mando, dias de
descanso desde a partida anterior, momento da temporada e regime de público.

A ausência de vazamento foi verificada experimentalmente. O procedimento de teste
altera o placar de uma partida situada no meio do conjunto, reconstrói todas as
variáveis explicativas e exige que nenhuma variável daquela mesma partida se
altere. Um segundo teste assegura que as variáveis das partidas subsequentes se
alterem, de modo a impedir que o primeiro seja satisfeito trivialmente.

A validação adotada foi temporal. Para prever determinada temporada, o conjunto
de treino compreende exclusivamente partidas anteriores a ela; dentro da
temporada avaliada, o histórico é ampliado a cada bloco de dez partidas e todos
os modelos são reajustados sobre ele. A uniformidade desse procedimento entre os
modelos é deliberada: caso apenas um deles reestimasse parâmetros ao longo da
temporada, a comparação mediria o protocolo de reajuste e não a qualidade dos
modelos. A imputação de valores ausentes utiliza estatísticas calculadas
exclusivamente sobre o conjunto de treino de cada dobra.

Foram avaliados cinco modelos: dois de referência, correspondentes à frequência
histórica dos resultados e à heurística de prever sempre a vitória do mandante; o
modelo de Poisson bivariado descrito na seção anterior, empregado como preditor;
uma regressão logística multinomial; e um modelo de \textit{gradient boosting}. As
métricas adotadas foram o \textit{log-loss} e o escore de Brier como principais,
a acurácia como métrica secundária e a curva de calibração como análise
complementar.
